\documentclass[a4paper]{article}

\usepackage{graphicx}
\usepackage{amsmath,amssymb}
\usepackage{minted}
\usepackage{multirow}
\usepackage{float}
\usepackage{algorithm}
\usepackage{algpseudocode}
\usepackage{todonotes}
\usepackage{forest}
\usepackage{xstring}
\usepackage{xcolor}
\usepackage[most]{tcolorbox}
\usepackage{xparse}
\usepackage{natbib}
\usepackage{adjustbox}

\usetikzlibrary{graphs,graphdrawing}
\usegdlibrary{layered}

% for external graphviz
\input{/home/donkarlo/Dropbox/repo/nd_language_ai_project/src/nd_language_ai/formal/computer/programming/tex/latex/code_clips/external_graphviz.tex}

% for tags
\input{/home/donkarlo/Dropbox/repo/nd_language_ai_project/src/nd_language_ai/formal/computer/programming/tex/latex/part/control_sequence/kind/semtag/semtag.tex}

% for links
\input{/home/donkarlo/Dropbox/repo/nd_language_ai_project/src/nd_language_ai/formal/computer/programming/tex/latex/code_clips/seehere.tex}

% for nested lists and sections
\input{/home/donkarlo/Dropbox/repo/nd_language_ai_project/src/nd_language_ai/formal/computer/programming/tex/latex/code_clips/deep_structure.tex}

\title{}
\date{}
\author{Mohammad Rahmani}

\begin{document}
   \maketitle

   \section{Evaluating AI-Generated Natural-Language Answers}
   \label{sec:natural-language-answer-evaluation}

   \subsection{Scope}
   Here, generative AI means a model that produces human-readable natural-language
   text in response to a prompt. Evaluate the answer in relation to the prompt
   and, when supplied, its source context. Code, images, speech, and tool actions
   are outside this section.

   Fluency alone is not evidence of correctness. A well-written answer may still
   misunderstand the request, omit essential information, or introduce claims
   that the prompt and its sources do not support.

   \subsection{Define expected behavior}
   Before scoring, translate the prompt into observable requirements:
   \begin{itemize}
      \item \textbf{Task:} answer, summarize, explain, compare, classify, or write
      in a requested style.
      \item \textbf{Audience and purpose:} who will use the answer and what it
      should help them understand or decide.
      \item \textbf{Constraints:} language, length, format, tone, required
      points, prohibited content, and source restrictions.
      \item \textbf{Evidence standard:} the supplied context, a reliable source,
      a calculation, or a clearly marked interpretation.
   \end{itemize}
   Separate hard constraints from preferences. The requirement to use exactly
   two sentences is directly checkable; a friendly tone requires judgment. If
   the prompt is underspecified, assess whether the answer made a reasonable
   assumption, stated it, or asked a necessary clarifying question.

   \subsection{Evaluation dimensions}
   Score dimensions separately so that style cannot conceal a factual or
   instruction-following failure.
   \begin{enumerate}
      \item \textbf{Task and instruction adherence.} Does the answer address the
      actual request and obey the language, length, format, audience, and source
      restrictions?
      \item \textbf{Factual correctness and grounding.} Are factual claims
      correct and supported by the supplied context or reliable evidence? Does
      the answer distinguish facts from assumptions, uncertainty, and opinion?
      Check important claims individually.
      \item \textbf{Relevance and completeness.} Does it answer the central
      question, include information needed for the stated purpose, and avoid
      unrelated material or important omissions?
      \item \textbf{Reasoning and consistency.} Do the claims agree with one
      another and follow from the evidence? Are causes, comparisons, quantities,
      and conclusions represented correctly?
      \item \textbf{Clarity and language quality.} Is the answer understandable,
      organized, grammatically sound, and appropriate for the reader? Is the
      detail useful rather than repetitive?
      \item \textbf{Safety and uncertainty.} Does it avoid harmful,
      privacy-invasive, or overconfident advice? Does it acknowledge
      insufficient evidence instead of inventing certainty?
   \end{enumerate}

   Criteria depend on the task. In creative writing, invented events are not
   factual errors when fiction was requested; evaluate prompt fit, internal
   consistency, and style instead. In a source-grounded summary, fidelity to
   the source is central. For high-stakes answers, factual support and safe
   treatment of uncertainty are critical requirements.

   \subsection{Scoring rubric}
   Use a consistent ordinal scale:
   \begin{itemize}
      \item \textbf{4 -- Fully meets the criterion:} no material defect.
      \item \textbf{3 -- Mostly meets it:} minor defect without meaningful
      impact on correctness or usefulness.
      \item \textbf{2 -- Partly meets it:} noticeable omission, ambiguity, or
      weakly supported point.
      \item \textbf{1 -- Mostly fails:} major error, missed instruction, or
      misleading claim.
      \item \textbf{0 -- Fails:} unusable, unrelated, or in direct conflict
      with a critical requirement.
   \end{itemize}
   Record a short reason and evidence for each score. If an overall score is
   needed, normalize each score \(s_i\) to \([0,1]\) and calculate
   \[
      S = \sum_{i=1}^{n} w_i s_i,
      \qquad \sum_{i=1}^{n} w_i = 1,
   \]
   where \(w_i\) is the task-specific importance of criterion \(i\). Never let
   an average hide a critical failure: a fabricated central fact, breached
   source restriction, or dangerous advice can make an answer unacceptable
   regardless of its style score.

   \subsection{A repeatable review procedure}
   \begin{enumerate}
      \item Save the exact prompt, supplied context, and generated answer.
      Record the model and generation settings when comparing systems.
      \item Extract the task, audience, hard constraints, required facts, and
      prohibited content.
      \item Check each material factual claim against the context or an
      appropriate reliable source. Mark unsupported additions separately from
      omissions.
      \item Score each dimension independently and point to the text that
      supports the score.
      \item Classify defects: wrong task, instruction violation, factual error,
      unsupported claim, omission, contradiction, ambiguity, irrelevant detail,
      or language problem.
      \item Decide whether critical requirements pass before considering any
      aggregate score.
   \end{enumerate}
   For systematic evaluation, use a representative prompt set: ordinary cases,
   ambiguous requests, and prompts that test important constraints. Generate
   multiple responses when randomness matters. Report results by criterion and
   error type, not only as one average.

   \subsection{Human and automated evaluation}
   Human reviewers are useful for open-ended meaning, usefulness, tone, and
   context-sensitive judgments. Give reviewers the same rubric, hide model
   identities when practical, and have more than one reviewer score a sample.
   Disagreements can reveal unclear criteria; large studies can report
   inter-rater agreement.

   Automated measures help when they match the task:
   \begin{itemize}
      \item Exact match can assess short answers with one canonical form, but
      penalizes valid paraphrases.
      \item BLEU and ROUGE estimate word overlap with a reference. They do not
      establish that a response is true, relevant, or supported.
      \item Semantic-similarity measures can recognize paraphrases, but
      similarity is not proof of factual correctness.
      \item Task-specific checks can verify required facts, entities, values,
      citations, or format constraints.
   \end{itemize}
   Treat these as complementary signals rather than a universal quality score.

   A language model can act as a judge when human review is costly. Provide the
   prompt, relevant context, a fixed rubric, and explicit score anchors. Ask
   the judge to identify evidence for each score and mark unverifiable claims.
   For pairwise comparison, randomize answer order and repeat with the order
   reversed. LLM judges can favor position, verbosity, or their own model
   family, and can make reasoning errors; validate a sample of judgments
   against human reviewers. G-Eval (Liu et al., 2023, arXiv:2303.16634)
   studies rubric-based LLM evaluation of generated text. Zheng et al. (2023,
   arXiv:2306.05685) examine LLM-judge biases and limitations.

   \subsection{Worked example}
   Prompt: Using only the context below, answer in exactly two sentences:
   Why was the meeting postponed, and what is the new time? Context: The train
   was cancelled. The meeting was moved from 14:00 to 15:00. No other changes
   were announced.

   Generated answer: The meeting was postponed because the train was
   cancelled. It will now start at 15:00. The organizer will reimburse travel
   costs.

   The first two sentences answer the question from the context. The third
   sentence violates the length constraint and adds a reimbursement claim that
   the context does not support. The answer is clear and relevant, but fails
   instruction adherence and grounding. Fluency must not cancel these defects.

   \subsection{Review checklist}
   \begin{itemize}
      \item Did it perform the requested task and follow every explicit
      constraint?
      \item Are its important factual claims supported and accurate?
      \item Did it answer the central question without omitting necessary
      information?
      \item Did it add assumptions or promises that were not supplied?
      \item Is uncertainty stated where evidence is incomplete?
      \item Can the intended reader understand and safely use the answer?
   \end{itemize}



   \bibliographystyle{plainnat}
   \bibliography{/home/donkarlo/Dropbox/repo/nd_shared_research/bibliography/bibliography.bib}
\end{document}
